% Generated draft: requires mathematical and source review.
\documentclass[11pt]{article}
\usepackage[margin=25mm]{geometry}
\usepackage{fontspec}
\setmainfont{Latin Modern Roman}
\usepackage{amsmath,amssymb,mathtools,unicode-math}
\setmathfont{Latin Modern Math}
\usepackage{enumitem,needspace,xurl,hyperref}
\setlength{\parindent}{0pt}
\setlength{\parskip}{5pt}
\setlength{\emergencystretch}{3em}
\newcommand{\N}{\mathbb{N}}
\newcommand{\Z}{\mathbb{Z}}
\newcommand{\Q}{\mathbb{Q}}
\newcommand{\R}{\mathbb{R}}
\newcommand{\C}{\mathbb{C}}
\newcommand{\F}{\mathbb{F}}
\newcommand{\A}{\mathbb{A}}
\newcommand{\PP}{\mathbb{P}}
\newcommand{\E}{\mathbb{E}}
\newcommand{\eps}{\varepsilon}
\newcommand{\dd}{\,\mathrm{d}}
\begin{document}
\section*{UnsolvedMath ID: 7600016}
\textbf{Source problem code:} catalogue entry
\subsection{Definitions and mathematical statement}
For a relational language \(L\), a structure \(M\) is ultrahomogeneous if every isomorphism between finite induced substructures extends to an automorphism of \(M\). A parameter-free definable reduct of \(M\) is a relational structure \(N\) on the same underlying set whose basic relations are first-order definable in \(M\) without parameters. Reducts are customarily counted up to parameter-free interdefinability: \(N_1\sim N_2\) if every basic relation of either is definable without parameters in the other.

Under these standard reduct conventions, the question asks whether every finitely presented homogeneous structure \(M\) in a finite relational language satisfies
\[
 \#\bigl(\{N:N\text{ is a parameter-free definable reduct of }M\}/\!\sim\bigr)<\infty.
\]
This is not the trivial question of deleting a subset of the finitely many relation symbols of \(L\); definable reducts may use relations not named in \(L\).

The excerpt does not define ``finitely presented'', nor explicitly specify countability, infinitude, the meaning of homogeneous, or the equivalence used to count reducts. In the usual countable homogeneous setting, one possible finite-constraint condition is finite boundedness: \(\operatorname{Age}(M)\) consists exactly of the finite \(L\)-structures avoiding induced copies of a fixed finite list of forbidden structures. That condition is not silently substituted for finite presentation here. The supplied summary calls the question a version of Thomas's conjecture, but does not resolve these formulation choices.
\subsection{Short English statement}
Must a finitely presented homogeneous structure in a finite relational language have only finitely many parameter-free definable reducts up to interdefinability? Finite presentation and the intended countability and reduct conventions need source definitions.
\subsection{Sources}
\begin{itemize}
\item Local UnsolvedMath registry, permanent ID 7600016, source code catalogue entry; update 2026-08-06; no proposer or proposal year supplied.
\item Wikipedia, \emph{List of unsolved problems in mathematics}, ``Model theory and formal languages'', bullet 16, pinned revision 1366636767: \url{https://en.wikipedia.org/w/index.php?title=List_of_unsolved_problems_in_mathematics&oldid=1366636767#Model_theory_and_formal_languages}. Registry access date 2026-07-29; MediaWiki extraction. Wikipedia text is CC BY-SA 4.0, requiring attribution and ShareAlike handling. Provenance says unpublished pending primary-source review, despite the registry published flag.
\item The supplied survey summary associates the question with Thomas's conjecture and mentions results for metrically homogeneous graphs and related structures. It supplies no exact primary citation or verified complete resolution.
\end{itemize}
\end{document}
